\begin{theorem}[Differential-form classifier symmetry obligation]
\label{thm:diffform-classifier-symmetry-obligation}
For rows carried by $\mathsf{DiffFormBHistCarrier}_{M,T}$, the classifier of
\autoref{def:diffform-bhist-form-classifier} is symmetric: if
$\omega\sim_{\DiffFormUp}\eta$, then $\eta\sim_{\DiffFormUp}\omega$.
\end{theorem}
\begin{proof}
Unpack the given classifier row through
\autoref{def:diffform-bhist-form-classifier}. The degree comparison and the
tensor, antisymmetry, and source-ledger comparisons are visible $\hsame$
relations on $\mathsf{BHist}$ rows, so symmetry of $\hsame$ reverses each of
those clauses.

The matched scalar endpoints are compared by the scalar-source classifier
retained from the $\NameCert_{\ManifoldUp}$ and $\NameCert_{\TensorProductUp}$
sources. Reversing that retained classifier comparison at each bundled
$\mathsf{ProbeBundle}$ row gives the scalar clause for
$\eta\sim_{\DiffFormUp}\omega$, and the reversed $\hsame$ ledger clauses supply
the remaining classifier fields.
\end{proof}

\begin{theorem}[Differential-form classifier transitivity obligation]
\label{thm:diffform-classifier-transitivity-obligation}
For rows carried by $\mathsf{DiffFormBHistCarrier}_{M,T}$, the classifier of
\autoref{def:diffform-bhist-form-classifier} is transitive: if
$\omega\sim_{\DiffFormUp}\eta$ and $\eta\sim_{\DiffFormUp}\theta$, then
$\omega\sim_{\DiffFormUp}\theta$.
\end{theorem}
\begin{proof}
Unpack the two classifier rows using
\autoref{def:diffform-bhist-form-classifier}. Transitivity of $\hsame$ composes
the degree comparisons, and the same transitive $\hsame$ composition applies
to the tensor, antisymmetry, and source-ledger $\mathsf{BHist}$ rows.

For each matched multilinear row in the common $\mathsf{ProbeBundle}$ surface,
compose the scalar endpoint comparisons through the scalar-source classifier
retained by $\NameCert_{\ManifoldUp}$ and $\NameCert_{\TensorProductUp}$. The
composed scalar clauses and the composed $\hsame$ clauses give precisely the
classifier fields required for $\omega\sim_{\DiffFormUp}\theta$.
\end{proof}

\begin{theorem}[Differential-form carrier hsame transport]
\label{thm:diffform-carrier-hsame-transport}
Let $\omega$ be a row accepted by
$\mathsf{DiffFormBHistCarrier}_{M,T}$, and let $\omega'$ be obtained by
replacing its six displayed histories with $\hsame$-equal histories while
keeping the same $\NameCert_{\ManifoldUp}$ and $\NameCert_{\TensorProductUp}$
sources. Then $\omega'$ is also accepted by
$\mathsf{DiffFormBHistCarrier}_{M,T}$.
\end{theorem}
\begin{proof}
Unpack the carrier row of \autoref{def:diffform-bhist-form-carrier}. The
degree condition is transported by $\hsame$ for $\mathsf{UnaryHistory}$.

The probe bundle is a closed $\mathsf{ProbeBundle}$ spine, so its bundled
tangent rows are transported pointwise through the displayed manifold source
classifier. The tensor row is then transported through the
$\TensorProductUp$ classifier, retaining the same $\Cont$ ledgers after
source, probe, and endpoint histories are replaced by $\hsame$-equal rows.

It remains to transport the scalar and antisymmetry rows. The scalar endpoint
uses the retained scalar classifier, and its pointwise transport follows from
classifier stability in the source surface. Each antisymmetry row is a visible
permutation ledger over the same probe bundle; transporting the two scalar
endpoints of the swap row gives the required signed comparison. Repacking
these transported fields gives the carried row for $\omega'$.
\end{proof}

\begin{theorem}[Differential-form classifier stability]
\label{thm:diffform-form-classifier-stability}
For rows carried by $\mathsf{DiffFormBHistCarrier}_{M,T}$, the classifier of
\autoref{def:diffform-bhist-form-classifier} is stable under carrier
transport: if $\omega\sim_{\DiffFormUp}\eta$ and both rows are transported to
$\omega'$ and $\eta'$ by componentwise $\hsame$ comparisons, then
$\omega'\sim_{\DiffFormUp}\eta'$.
\end{theorem}
\begin{proof}
The degree comparison is the transitive composite of the transport comparison
for $\omega$, the original degree comparison, and the symmetric transport
comparison for $\eta$.

For each bundled multilinear probe row, use the tensor-source classifier to
identify the matched tensor rows after transport. The scalar endpoints are
then compared by the scalar classifier stability already retained in the
upstream scalar surface.

The antisymmetry and source-ledger clauses are pure $\hsame$ clauses on
visible BHist rows. Composing their original comparisons with the two
transport comparisons gives the required ledger matches for the transported
rows.
\end{proof}

\begin{theorem}[Differential-form probe-slot arity transport]
\label{thm:diffform-probe-slot-arity-transport}
Let
$$
\omega=(d,\mathcal{P},\tau,\sigma,\alpha,\lambda)
$$
and
$$
\omega'=(d',\mathcal{P}',\tau',\sigma',\alpha',\lambda')
$$
be rows carried by $\mathsf{DiffFormBHistCarrier}_{M,T}$. Suppose
$\hsame(d,d')$ and the displayed $\mathsf{ProbeBundle}$ spines of
$\mathcal{P}$ and $\mathcal{P}'$ are matched componentwise through the
tangent-source classifier supplied by $\NameCert_{\ManifoldUp}$. Then the
visible arity of the probe slots in $\mathcal{P}$ is the same as the visible
arity of the probe slots in $\mathcal{P}'$. Moreover, every transported
multilinear tensor row remains classified by the $\TensorProductUp$ source
over those transported tangent probes.
\end{theorem}
\begin{proof}
Unpack the carrier rows in
\autoref{def:diffform-bhist-form-carrier}. The degree entries are unary
history rows, so $\hsame(d,d')$ transports the displayed degree without
creating or deleting a probe slot.

The probe component is a closed $\mathsf{ProbeBundle}$ spine. Induct on this
spine. The empty spine has arity zero on both sides. In the cons case, the
head tangent probe is transported through the classifier retained by
$\NameCert_{\ManifoldUp}$, and the induction hypothesis preserves the arity of
the tail. Therefore the visible slot count is preserved at each spine
constructor.

It remains to place the transported multilinear row back in the tensor source.
The tensor component of a carried differential-form row is accepted only as a
$\TensorProductUp$ row over the bundled tangent probes. Since the probe
classifier transport is componentwise and the tensor row retains its displayed
$\Cont$ provenance, tensor-source classifier transport reclassifies $\tau'$ as
a multilinear row over the transported probe bundle.
\end{proof}

\begin{theorem}[Differential-form adjacent-swap involution ledger]
\label{thm:diffform-adjacent-swap-involution-ledger}
Let
$$
\omega=(d,\mathcal{P},\tau,\sigma,\alpha,\lambda)
$$
be carried by $\mathsf{DiffFormBHistCarrier}_{M,T}$. If a displayed row of the
antisymmetry ledger $\alpha$ records an adjacent swap of two neighboring probe
slots and classifies the resulting scalar endpoint with the signed endpoint
required by antisymmetry, then $\alpha$ also determines the reverse adjacent
swap row. Applying the displayed adjacent swap and then its reverse returns
endpoints classified with the original scalar endpoint under the retained
scalar classifier, and the provenance for both swap rows remains recorded in
$\lambda$.
\end{theorem}
\begin{proof}
Read the chosen adjacent-swap row from the visible antisymmetry ledger of
\autoref{def:diffform-bhist-form-carrier}. Its endpoints are scalar rows
classified by the scalar classifier retained from the manifold and tensor
sources, and its provenance is one of the rows listed in $\lambda$.

The reverse row uses the same two neighboring probe slots in the opposite
order. The scalar comparison for the reverse row is obtained by symmetry of
the retained scalar classifier, while the comparison between the double-swap
endpoint and the original endpoint is the transitive composite of the forward
and reverse scalar comparisons.

No hidden row is introduced in this construction. Both the displayed swap row
and the reverse swap row are retained as BHist provenance in $\lambda$, so the
two applications return to endpoints classified with the original scalar
endpoint inside the same carried differential-form row.
\end{proof}

\begin{theorem}[Differential-form wedge degree additivity ledger]
